
\item Suppose you are given a magic black box that somehow answers the following
decision problem \emph{in polynomial time}: given positive integers $n$ and $k$,
does $n$ have a prime factor less than $k$?

Formally, a \EMPH{prime factorization} of a positive integer $n$ is a multiset
of prime numbers $\{p_1, p_2, \dots, p_m\}$ such that $n = p_1 p_2 \cdots p_m$.
For example, the prime factorization of $12$ is $\{2, 2, 3\}$, and the prime
factorization of $13$ is $\{13\}$.

Suppose you are given a magic black box that somehow answers the following
decision problem \emph{in polynomial time}:
\begin{itemize}
\item \mathsc{Input:} Positive integers $n$ and $k$.
\item \mathsc{Output:} \mathsc{True} if $n$ has a prime factor less than $k$,
and \mathsc{False} otherwise.
\end{itemize}
Using this black box as a subroutine, describe an algorithm that solves the
following \EMPH{search problem} \emph{in polynomial time}:
\begin{itemize}
\item \mathsc{Input:} A positive integer $n$.
\item \mathsc{Output:} A prime factorization of $n$ (i.e., a multiset of prime
numbers $\{p_1, p_2, \dots, p_m\}$ such that $n = p_1 p_2 \cdots p_m$).
\end{itemize}
\Hint{You can use the magic box more than once.  The input to the magic box is a pair of positive integers $(n, k)$.}

\emph{Unlike in the problem session solutions}, you do not need to argue the
correctness of your algorithms. But you still do need to argue that they run in
polynomial time.

\begin{Solution}
We can find the prime factorization of $n$ using the following algorithm:
\begin{algo}
\textul{$\mathsc{FindPrimeFactors}(n)$:}\+
\\	$P \gets \emptyset$ \quad\text{(the multiset of prime factors found so far)}\+
\\	$k \gets 2$\+
\\	\textul{while} $n > 1$ \textul{do}\+
\\		\textul{if} $\mathsc{DecideHasSmallPrimeFactor}(n, k)$ \textul{then}\+
\\			$P \gets P \cup \{k\}$\+
\\			$n \gets n / k$\-
\\		\textul{else}\+
\\			$k \gets k + 1$\-
\end{algo}
\end{Solution}